% ===== 原文 PDF pp.17--19（印刷 pp.16--18）=====
% tr-01 · tr01c.tex · tag 区间 2.15--2.27（13 条）
% 内容：N=1 超空间形式 (2.15)--(2.20) / SYM 拉格朗日量 (2.21)--(2.25) / §2.3 标题 / 标量势与模空间 (2.26)--(2.27) / 图 2
% 边界：本文件写到 PDF p.19 末（含 p.20 首句「and equivalently D=F=0」及其脚注 11）

% [原文 p.16]
对我们的后续而言，一个重要的 $\cN=1$ 超场是所谓的\textbf{实矢量超场} $V$，它把矢量多重态中的 Weyl 费米子 $\lambda_\alpha$ 与矢量场 $A_\mu$ 重新打包成紧凑的形式\footnote{这里我们取 Wess--Zumino 规范。}
\begin{equation}
  V=-i\theta\sigma^\mu\bar\theta A_\mu+i\theta\theta\bar\theta\bar\lambda-i\bar\theta\bar\theta\theta\lambda+\frac12\theta\theta\bar\theta\bar\theta D\,,
  \tag{2.15}
\end{equation}
其中 $D$ 是辅助场。另一个在 4d $\cN=1$ 手征多重态中包含物质场的超场称为\textbf{手征超场} $Q$
\begin{equation}
  Q^A=q^A+i\psi^A\theta+F^A\theta\theta\,,
  \tag{2.16}
\end{equation}
它在 $\cN=1$ 的 $\mathrm{U}(1)_r$ 下带荷，记为 $r_{\cN=1}$。同样，$A$ 是味指标，$F^A$ 是辅助场，而它在超对称下按全导数变换。
% [Tutor Note: (2.15) 是 Wess--Zumino 规范下的展开：由于 (2.15) 中只出现 $\theta\sigma^\mu\bar\theta A_\mu$、$\theta\theta\bar\theta\bar\lambda$、$\bar\theta\bar\theta\theta\lambda$、$\theta\theta\bar\theta\bar\theta D$ 四项，$V$ 的其余分量（如 $\theta\theta$、$\bar\theta\bar\theta$ 分量）已被规范变换消去。注意 $V$ 中并不出现场强 $F_{\mu\nu}$，这正是后文要把矢量多重态改写为规范微子超场 $W_\alpha$（见 (2.21)）的原因。]

然后，$\cN=2$ 矢量多重态可以表示为 $\cN=1$ 矢量超场与 $\cN=1$ 手征超场的组合：
\begin{equation}
  V_{\cN=2}\to(V,Q)\,,
  \tag{2.17}
\end{equation}
其中手征超场 $Q$ 在规范群 $G$ 的伴随表示中，且 $r_{\cN=1}=\frac23$。这里我们用了 $\cN=1$ 与 $\cN=2$ 超共形 R 对称性之间的关系
\begin{equation}
  r_{\cN=1}=\frac23 r_{\cN=2}+\frac43 I_3\,,
  \tag{2.18}
\end{equation}
其中 $I_3$ 是 $\SU(2)_R$ 的 Cartan 生成元。另一方面，$\cN=2$ 超多重态可以由两个 $\cN=1$ 手征超场表示：
\begin{equation}
  H_{\cN=2}\to(Q,\tilde Q)\,,
  \tag{2.19}
\end{equation}
同样有 $r_{\cN=1}=\frac23$。
% [Tutor Note: 用 (2.18) 核对 $\Phi$ 的 $r_{\cN=1}$：由 (2.13) 后的讨论，$\Phi$ 在 $\cN=2$ 下带 $r_{\cN=2}=1$、且它是 $\SU(2)_R$ 单态故 $I_3=0$，代入 (2.18) 得 $r_{\cN=1}=\frac23\cdot1=\frac23$，与原文一致。$\frac43 I_3$ 这一项的存在说明 $\cN=2$ 的 $\fu(1)_r$ 与 $\SU(2)_R$ 的 Cartan 在 $\cN=1$ 的语言里混合成了同一个 $\fu(1)$；只有 $\frac23 r_{\cN=2}$ 而无 $\frac43 I_3$ 时，$\mathcal N=2$ 的 R 对称性无法正确嵌入。]

最一般的 4d $\cN=1$ 超对称拉格朗日量可以由手征超场的超势以及可能涉及手征超场与实矢量超场 $V$ 之间耦合的 Kähler 势构造出来。通过限制场的内容，并对超势与 Kähler 势作特殊的选择，拉格朗日量可以具有增强的 $\cN=2$ 超对称性。示意地，这样的拉格朗日量具有如下形式
\begin{equation}
  \cL_{\cN=2}=\cL^{\mathrm{SYM}}_{\cN=2}+\cL^{\mathrm{matter}}_{\cN=2}\,.
  \tag{2.20}
\end{equation}

% [原文 p.17]
其中第一项包含 4d $\cN=2$ 矢量多重态的贡献，第二项包含 $\cN=2$ 超多重态的贡献（超多重态可能与矢量多重态相耦合）。

为了写出矢量多重态的规范不变作用量、使熟悉的 Yang--Mills 作用量 $F_{\mu\nu}F^{\mu\nu}$ 超对称化，把矢量多重态 $V$ 改写为所谓的\textbf{规范微子超场} $W_\alpha$ 会比较方便，因为 $V$ 并不直接包含规范场强 $F_{\mu\nu}$；$W_\alpha$ 具有如下展开：
\begin{equation}
  W_\alpha=-i\lambda_\alpha+iF_{\mu\nu}(\sigma^{\mu\nu}\theta)_\alpha+D\theta_\alpha+\theta\theta(\sigma^\mu\partial_\mu\bar\lambda_\alpha)\,,
  \tag{2.21}
\end{equation}
其最低分量 $\lambda_\alpha$ 是 $r_{\cN=1}=1$ 的 Weyl 规范微子。它与 $V$ 通过超协变导数 $D_\alpha\equiv\frac{\partial}{\partial\theta^\alpha}+i\sigma^\mu_{\alpha\dot\beta}\bar\theta^{\dot\beta}\partial_\mu$ 相联系：
\begin{equation}
  W_\alpha=-\frac14\bar D^2D_\alpha V\,.
  \tag{2.22}
\end{equation}
% [Tutor Note: 对照 (2.15) 与 (2.21)：$A_\mu$ 出现在 $V$ 中而 $F_{\mu\nu}$ 只出现在 $W_\alpha$ 中，因此不能直接用 $V$ 写规范不变的 $\int F_{\mu\nu}F^{\mu\nu}$；这正是引入 $W_\alpha$ 的动机。另注意 (2.22) 的 $\bar D^2D_\alpha$ 中 $D_\alpha$ 是超协变导数（不是辅助场 $D$），二者只是同名，读者勿混。]

现在，由 (2.17) 而来的 $\cN=2$ 矢量多重态分裂为一个 $\cN=1$ 规范微子超场 $W_\alpha$ 与一个 $\cN=1$ 手征超场 $\Phi$。相应的 $\cN=2$ 超 Yang--Mills 拉格朗日量取如下超空间积分的形式
\begin{equation}
  \cL^{\mathrm{SYM}}_{\cN=2}=\frac{1}{8\pi i}\int d^2\theta\,\Tr\left(\tau W_\alpha W^\alpha+\mathrm{c.c.}\right)
  +\frac{\Im(\tau)}{4\pi}\int d^4\theta\,\Tr\left(\Phi^\dagger \ee^{\mathrm{adj}(V)}\Phi\right).
  \tag{2.23}
\end{equation}
这里第一项编码了 $\cN=1$ SYM 的动能项，而第二项表示手征超场在与伴随表示中的矢量超场作最小耦合下的 $\cN=1$ 动能项。复化规范耦合 $\tau$ 同时包含 Yang--Mills 耦合与 theta 角，
\begin{equation}
  \tau=\frac{\theta_{\mathrm{YM}}}{2\pi}+\frac{4\pi i}{g^2}\,.
  \tag{2.24}
\end{equation}
% [Tutor Note: (2.23) 第二项的前因子用 (2.24) 化简易验：$\Im(\tau)=\Im\left(\frac{\theta_{\rm YM}}{2\pi}+\frac{4\pi i}{g^2}\right)=\frac{4\pi}{g^2}$，故 $\frac{\Im(\tau)}{4\pi}=\frac{1}{g^2}$，即手征超场动能项的标准系数；而第一项 $\frac{1}{8\pi i}\tau W_\alpha W^\alpha$ 中的 $\theta_{\rm YM}$ 部分给出拓扑项、$\frac{1}{g^2}$ 部分给出 $-\frac{1}{4g^2}F_{\mu\nu}F^{\mu\nu}$。两相配合正是 $\cN=2$ 的要求。]
(2.23) 中的相对前因子由 $\SU(2)_R$ 对称性固定，它保证整个拉格朗日量保持 $\cN=2$ 超对称。超多重态的 $\cN=2$ 拉格朗日量包含 (2.19) 中 $(Q,\tilde Q)$ 在与规范群 $G$ 的某个表示 $\rho$ 中的实矢量超场 $V$ 作最小耦合下的动能项，即\footnote{在一般情况下，还可以为超多重态添加一个质量项 $\int d^2\theta\,m\tilde QQ$。为简单起见，我们在后续讨论中省略这一项。}
\begin{equation}
  \cL^{\mathrm{matter}}_{\cN=2}
  =\int d^4\theta\left[Q^\dagger \ee^{\rho(V)}Q+\tilde Q^\dagger \ee^{\bar\rho(V)}\tilde Q\right]
  +\int d^2\theta\,\tilde Q\rho(\Phi)Q+\mathrm{c.c.}
  \tag{2.25}
\end{equation}

% [原文 p.18]
第二项写成一个 $\cN=1$ 超势，\footnote{回忆一下，$\cN=1$ 超场 $Q,\tilde Q,\Phi$ 都具有 $r_{\cN=1}=\frac23$。}它编码了 $\cN=2$ 矢量多重态与超多重态之间相互作用的一部分，而该耦合由 $\SU(2)_R$ 对称性固定为 $1$。注意由上述结果可见，一般的无质量 $\cN=2$ 拉格朗日量中除复化规范耦合 $\tau$ 之外没有任何参数。

\SUB{2.3}{超对称真空模空间}

我们在上一节构造了包含矢量多重态与超多重态的最一般的 $\cN=2$ 拉格朗日量，而通过这个拉格朗日量可以研究的一件具体事情是真空的（经典）\textbf{模空间}。它将来自标量势 $V(\Phi,q,\tilde q)$，后者由拉格朗日量 (2.20) 在积掉矢量超场中的辅助场 $D$ 与手征超场中的 $F$ 之后得到。标量势 $V$ 是相应于所谓 D 项与 F 项的平方和 [117]，
\begin{equation}
  V(\Phi,q,\tilde q)=\frac12\Tr\left(D^2\right)+F\bar F\,,
  \tag{2.26}
\end{equation}
其中辅助场 $D$ 与 $F$ 由物理标量场 $\Phi,q,\tilde q$ 确定，这一点我们将在下文看到。

\begin{center}
  \includegraphics[width=0.72\linewidth]{fig2.png}
\end{center}
\figtitle{2}{一般 $\cN=2$ 理论的真空模空间。}
% [Tutor Note: 图 2 中三个分支的交点标为 SCFT：在原点处所有标量 VEV 为零、共形对称性未破缺，理论才是 SCFT；$\cM_{\mathrm{CB}}$（库仑分支）是图中那张水平面（黑线），$\cM_{\mathrm{HB}}$（希格斯分支）是蓝色椭圆，$\cM_{\mathrm{MB}}$（混合分支）是红色曲面/红线包出的区域。三者在原点合并，正是 (2.27) 中「零点集」的几何图像。]

超对称真空的模空间是标量势的零点集，
\begin{equation}
  \cM^{\mathrm{classical}}=\{\Phi,q,\tilde q\,|\,V(\Phi,q,\tilde q)=0\}\,,
  \tag{2.27}
\end{equation}
% [Tutor Note: 为什么零点集等价于 $D=F=0$：$V$ 是平方和（$\frac12\Tr D^2+F\bar F$，两项各自非负），故 $V=0$ 当且仅当每一项都为零；而辅助场没有动能项，其「运动方程」就是代数方程 $D=F=0$。这正是原文在 p.20 首句所说的「and equivalently $D(\Phi,q,\tilde q)=F(\Phi,q,\tilde q)=0$」；脚注 11 补充了另一种等价视角——要求费米子场的超对称变分为零。]
以及等价的 $D(\Phi,q,\tilde q)=F(\Phi,q,\tilde q)=0$。\footnote{辅助场的消失条件也可以由要求费米子场的超对称变分消失而得到，而后者是该位形为超对称所必需的。}
% [边界] 本句已写完，下一位 tr-02 从 PDF p.20 第 2 句「To be concrete, in the following let us consider SQCD ...」续
